\chapter{The constant increase of entropy is the basic law of the universe (I Was Stupid, So Stupid)}
The eighth episode of Puella Magi Madoka Magica is an example of a hazard typical of shows which, like Madoka Magica, rely heavily on the "shocking swerve" \cite{ref54} to keep the plot moving: namely, an excellent character episode ends up largely overshadowed by a devastating plot development in its final moments. As such, we're going to largely ignore that plot development until next chapter, and focus instead on the first 18 minutes or so of the episode.

Those first 18 minutes comprise one of the most focused episodes of the show; we spend all but a couple of scenes with Sayaka. Even one of those scenes is between Sayaka's satellite characters Kyousuke and Hitomi, and while it starts apart from her, its continuation is shown from her perspective. Almost all of the episode is thus dedicated to showing Sayaka's final breakdown, as the stress of being a magical girl, her complete loss of self-esteem after the revelation of what Kyubey has done to her body, and her refusal to purify her Soul Gem combine to utterly destroy her mental state.

This breakdown makes clear not only what has been going on with Sayaka, but a major component of what being a magical girl signifies in this world: to be a magical girl is to struggle constantly with depression—and the portrayal the show gives of depression is remarkably accurate to the lived experience of the disorder. The darkness which accumulates in each girl's Soul Gem as she uses magic is an external force that corrupts and distorts her worldview and thought processes, generating a cycle of despair that only darkens the gem faster and farther. Events in her life do not cause this cycle, but they can exacerbate and accelerate it \cite{ref55}. 

The common thread among all the magical girls we have seen is depression, though each expresses this state differently. Kyoko, for example, shows few outward signs of depression at first glance, but the previous episode made quite clear that she has totally cut herself off from humanity, burying herself in self-indulgence of her hedonistic impulses in order to escape the pain of her past. The images used in her flashback—in which she is frequently slumped over alone in a large space, or the only colored doll amongst a field of monochrome cut-outs—make clear the isolation she feels. Further, as someone who has experienced true hunger, she strongly associates food with comfort, as witness her offer of snacks to Sayaka shortly before the latter becomes a witch in this episode, to Madoka when they hatch their plan to try to retrieve Sayaka from her witch-state, and the piles of wrappers and containers in the apartment where she stashes Sayaka's body, both in the next episode. Food, for Kyoko, is a way to experience a connection to other people and the world that she consciously denies herself, as well as a means of self-medicating her feelings. In other words, the reason she is constantly eating may be because she has an eating disorder, one which could cause her serious harm if she were still fully human.

At the other end of the spectrum from Kyoko is Sayaka, who expresses her depression outwardly. In Sayaka's case, the frequent water imagery associated with her now carries a double meaning; she is, as her witch-form at the end of the episode suggests, playing the titular role of The Little Mermaid, referring here to the Hans Christian Anderson story about the maiden who sacrifices her physical form to be with the man she loves, only to have him (entirely unaware of her sacrifice) choose another \cite{ref56}, as opposed to the now-familiar uplifting Disney movie about the otherkin who, through much hardship and sacrifice, is able to move into the species she believes she belongs in \cite{ref57}. However, the imagery is equally a suggestion that Sayaka feels she is drowning, overwhelmed by a series of devastating losses and disappointments. At the start of this episode, she has already faced the disappointment of Kyousuke basically ignoring her from the moment he leaves the hospital, and of discovering that she's not the shining paragon of justice she hoped to become, but rather weaker than most magical girls. Compounding this, her depression distorts her perception, as depression tends to do, causing her to interpret everything that happens to her and that she does in the worst light possible. Thus, in her eyes, the changes to her physical form represent a total loss of humanity and strip her of all worth, Madoka's concern for her friend becomes reinterpreted as pity, and the misogynists on the train become representatives of a world not worth fighting to save anymore.

Central to Sayaka's collapse is her wish to heal Kyousuke, and the resulting lack of a relationship with him. It's an important choice that we do not hear his conversation with Hitomi; although the last episode and Rebellion both suggest that they do indeed start dating, Sayaka has no way of knowing this; she assumes that it's so both because she is putting the worst possible interpretation on everything that happens, and because her intense sense of inferiority leads her to elevate those around her. It is unthinkable to Sayaka that Hitomi won't get what she wants, because in Sayaka's distorted thinking Sayakas never get what they want and non-Sayakas always do. Sayaka's outburst on the train is likewise a response to her feelings about herself and Kyousuke; the talk of men callously abandoning women who care about them hits far too close to home for Sayaka not to lash out, and in (probably) murdering those men she is equally expressing her rage at Kyousuke, hence her comment shortly after that she no longer remembers what it was that she thought was worth fighting for. 

Very different from Sayaka is the other collapse we see in this episode, Homura's. Like Sayaka, she is pushing against exhaustion. There have been moments throughout the series where we see quick flashes of Homura expressing anger or grief—and occasionally with other characters as well, such as during Sayaka's conversation with Madoka in the rain shelter. These scenes are somewhat ambiguous; they are clearly glimpses of the emotions Homura is trying to hide, but it is not entirely clear whether they are purely extradiegetic images for the audience's benefit or, given the revelation of Homura's true nature in this episode, glimpses of how certain scenes played out slightly differently in other timelines. In either case, her desperation and misery immediately after shooting Kyubey, when she pleads with Madoka to stop sacrificing herself and value her own life, is the first time we see her true face within the diegetic space of the show's depiction of this particular timeline. Even Madoka recognizes the significance of this moment, as it is the first time she fully recognizes that she has met Homura before.

The cool, efficient Homura we see the rest of the time is a portrait of "functional depression" or dysthymia, a less severe, but longer-lasting, form of depression than the more familiar major depressive disorder \cite{ref58}. Like Kyoko, she is in a sense self-medicating, but in her case her "drug" of choice is the relentless pursuit of her cause. She has no hope, no belief that she can or will succeed or that anything will ever get better, but she is able to keep going because there are things she has to do. So long as she can keep moving, keep working toward this unachievable goal, she never has to think about how hopeless her situation is or how much suffering she has endured. So long as she keeps working, she does not need to think about how isolated she is, or that by working to save Madoka she has isolated herself from Madoka. As Sayaka astutely observes, Homura has given up on life just as much as Sayaka has; the difference is that Homura still has something she believes in. So long as Homura can continue to convince herself that nothing matters except helping Madoka, she can convince herself that her own feelings of despair and isolation do not matter, and thus keep herself from becoming a witch. In this sense Homura continues to be Mami's opposite number, since Mami likewise kept herself going through the isolation and loneliness of being a magical girl by dedicating herself to a cause.

Even Madoka herself shows hints of a generally depressed attitude, in particular a low self-worth and, as Homura points out, readiness to sacrifice herself for the sake of others (that is, valuing others above herself consistently). This bears little resemblance to the Madoka we will eventually see in other timelines, and suggests that, much as bits and pieces of other time lines bleed through Madoka's perception by way of dreams and quick flashes, so too has a little bit of the depression inherent in being a magical girl. The incompleteness and inconsistency of this depiction can thus be read as a consequence of the fact that she isn't actually a magical girl on this timeline.

These highly accurate and varied depictions of depression as lived from the inside are quite suggestive of the mental health history of the creators of the show, particularly Gen Urobuchi in light of his afterword to Fate/Zero Volume 1, as I quoted in Chapters 4 and 5. Of course it is both unwise and extremely impolite to try to diagnose an artist through their work; it is critically important to distinguish between the implied author of a work, which is a construct created by the audience looking through the work in an attempt to see the person on the other side, and the actual author, who is a full and real person utterly unknowable to anyone except the people who actually know them. Indeed, in this case the actual authors are a multitude of people, including writers, directors, character designers, animators, and voice actors, while the implied author remains as always an individual. That said, the implied Urobuchi seems pretty depressed, in a state of despair sufficient that he cannot even imagine a positive world where good things happen or create art depicting such a world.

In the remaining episodes, it becomes clear that Kyubey's goal is to use the energies released by the magical girls and witches to combat entropy, holding back the inevitable decay of the universe by, in effect, transferring that decay from physical reality to the emotional state of the magical girls. The universe in which these characters live is suffused with decay, and that decay is equivalent to depression and despair. But that decaying universe itself exists within the mind of a depressed implied author straining to transcend his limitations; Kyubey's process, in other words, is an attempt to stave off and exorcize depression by transferring it from the mind of his author to the minds of the characters.

Yet this is also the episode in which the other characters, with the exception of Madoka, begin unambiguously treating Kyubey as an enemy. The end of the episode in particular makes it very clear, in framing, music, and timing, that Kyubey is a complete monster, smugly profiting from the torture of magical girls who had no idea what they were agreeing to. Even there, however, Kyubey's attempted manipulation of Madoka into wishing to save Sayaka follows immediately on the scene with the misogynists on the train, who talk about their plans to manipulate the women in their lives, whom they regard as inferior creatures. Sayaka's statement shortly after the Kyubey scene, that she no longer remembers what she was trying to save, thus embeds Kyubey's monstrosity within the question of whether humanity is really any better. 

If Kyubey is the guardian and maintainer of the implied author's process of coping with depression, however, that paints the author in equally negative light, and further suggests that he is fully aware of and deliberately calling attention to how problematic this is. That, in turn, suggests that he is seeking another solution.

But to find it requires breaking the system he set up to begin with, which is to say the entire Madoka Magica universe. And to do that, he must first destroy the genre to which it belongs.
